\section{Expected Outcomes and Discussion}
The study will test whether personalized pace prediction reduces trip-time error, whether uncertainty-aware queries identify time-dependent constraints earlier, and whether deterministic verification reduces infeasible plans compared with language-model-only planning. No experimental results are claimed at this stage.
% TODO: Insert measured results, confidence intervals, ablation tables, and case studies after evaluation.
\section{Limitations}
The dataset represents a single rider, limiting generalization. Historical machine-readable road closure records are incomplete. Forecast errors, uncertain fuel consumption, and GPX measurement noise may affect recommendations. Civil twilight and P90 are planning criteria, not guarantees of safety. The initial scope is pre-trip decision support; live tracking is a potential extension.
\section{Conclusion}
This paper presents the design of Moto Touring Copilot, combining personalized travel-time prediction, time-aware environmental data integration, and deterministic verification. Future work will implement and evaluate the framework using the proposed experiments and document reproducible cases and decision traces.
